\documentclass[11pt,a4paper]{amsart}
\usepackage[margin=30mm]{geometry}
\usepackage{amsmath,amssymb,amsthm,mathtools}
\usepackage{enumitem,booktabs,longtable}
\usepackage[hidelinks]{hyperref}
\usepackage{url}
\makeatletter\g@addto@macro\UrlBreaks{\do\_\do\a\do\e\do\i\do\o\do\u}\makeatother
\newcommand{\lean}[1]{\nolinkurl{#1}}

\newcommand{\Q}{\mathbb{Q}}
\newcommand{\Z}{\mathbb{Z}}
\newcommand{\F}{\mathbb{F}}
\newcommand{\Qp}{\mathbb{Q}_p}
\newcommand{\Zp}{\mathbb{Z}_p}
\newcommand{\Fp}{\mathbb{F}_p}
\newcommand{\Reg}{\mathrm{Reg}}
\newcommand{\Et}{\widetilde{E}}
\newcommand{\Lam}{\Lambda}
\newcommand{\qp}{q_p}
\newcommand{\cL}{\mathcal{L}}
\newcommand{\Mb}{\overline{M}}

\theoremstyle{plain}
\newtheorem{theorem}{Theorem}[section]
\newtheorem{lemma}[theorem]{Lemma}
\newtheorem{proposition}[theorem]{Proposition}
\newtheorem{corollary}[theorem]{Corollary}
\theoremstyle{definition}
\newtheorem{definition}[theorem]{Definition}
\theoremstyle{remark}
\newtheorem{remark}[theorem]{Remark}

\begin{document}

\title[The height matrix modulo $p$ at split anomalous primes]{The height matrix modulo $p$ and the $p$-adic digits of the height at split anomalous primes}
\author{Yoshikazu Hosomura}
\address{Independent researcher, Japan}
\email{wvbks0@gmail.com}
\subjclass[2020]{Primary 11G50; Secondary 11G07, 11S31, 11Y40}
\keywords{$p$-adic height, $p$-adic regulator, anomalous prime, $p$-adic sigma function, division polynomial, Fermat quotient}

\begin{abstract}
Let $E/\Q$ be an elliptic curve and $p\ge5$ a prime of good ordinary reduction with $\#\Et(\Fp)=p$ and $E(\Qp)[p]\ne0$. In a previous paper we showed that, under mild hypotheses, the cyclotomic $p$-adic height matrix of $E$ is then $p$-integral and has rank at most two modulo $p$. Here we determine it further. Modulo $p$ the matrix is given by two local constants and two homomorphisms $\chi,\gamma$ defined by the first two layers of the filtration of $E(\Qp)$. The second $p$-adic digit of the height involves the Fermat quotient of the denominator of the point, so the description in terms of the filtration does not continue. We show that all digits are truncations of a single identity involving two Iwasawa logarithms and an even power series in the formal logarithm whose coefficients are explicit invariants of $E$; this determines at which digit each invariant appears, and it shows that $\min(v_p(\Reg_p),k)$ in rank one is a function of $k$ layers of the filtration and $k-1$ digits of one Fermat quotient, and of no smaller set of data. A non-degeneracy condition used in these results, $v_p(\phi_p'(x(T)))=1$ at the points of order $p$, is proved for all $p\ge5$. Finally, a point deep in the formal filtration has large height; this gives an archimedean sufficient condition for the regulator to be a unit in rank two, which accounts for the observed frequency of non-unit regulators in our computations.
\end{abstract}

\maketitle

\section{Introduction}\label{sec:intro}

Let $E/\Q$ be an elliptic curve given by a global minimal Weierstrass equation, $p\ge5$ a prime of good ordinary reduction, and $\Lam=E(\Q)/E(\Q)_{\mathrm{tors}}$, of rank $r$. Let $h_p$ be the cyclotomic $p$-adic height and $\Reg_p$ the corresponding regulator. Suppose that $p$ is anomalous, $\#\Et(\Fp)=p$, that $p\nmid n:=\prod_\ell c_\ell$, and that ($\iota$) the reduction map $\Lam\to\Et(\Fp)$ is surjective. In \cite{P1} we proved that the height matrix $M$ is $p$-integral if and only if $E$ is \emph{split at $p$}, that is, $E(\Qp)[p]\ne0$, and that in the split case $\Mb=M\bmod p$ vanishes on $\ker\chi\times\ker\chi$, where $\chi\colon\Lam\to\Fp$ is the discrete logarithm of the reduction. This paper studies the split case in more detail.

\subsection{The matrix modulo $p$}
In the split case $E(\Qp)=D\oplus E_1(\Qp)$ with $D=E(\Qp)[p]$. Writing a point as $\chi(Q)T_0+R$ with $T_0\in D$ and $R\in E_1(\Qp)$, put $\gamma(Q)=t(R)/p\bmod p$. Then $\chi$ and $\gamma$ are homomorphisms $E(\Zp)\to\Fp$, and $\ker\chi\cap\ker\gamma=E_2(\Qp)$.

\begin{theorem}[Theorem~\ref{thm:K}]
There are constants $\lambda_{T,1},B_0\in\Fp$, depending only on $E/\Qp$ and on the choice of $T_0$, with $B_0\ne0$, such that for all $P,R\in\Lam$
\[
\Mb(P,R)=2\lambda_{T,1}\,\chi(P)\chi(R)+B_0\bigl(\chi(P)\gamma(R)+\gamma(P)\chi(R)\bigr).
\]
\end{theorem}

So $\Mb$ has rank two exactly when $(\chi,\gamma)\colon\Lam\to\Fp^2$ is surjective, and in rank $r\ge3$ the regulator has valuation exactly $r-2$ if and only if $p^{-1}h_p$ is non-degenerate modulo $p$ on $\Lam\cap E_2(\Qp)$ (Corollary~\ref{cor:K}).

\subsection{The second digit and beyond}
One might expect $v_p(\Reg_p)\ge2$ to be governed in the same way by the image of $\Lam$ in $E(\Qp)/E_3(\Qp)$. This is false. For $P$ with $\chi(P)\ne0$, write $x(nP)=a/d'^2$ in lowest terms and let $\Theta$ be the second Fermat quotient of $\phi_p(x(nP))$ (Section~\ref{sec:second}). Then
\[
n^2h_p(P)\equiv(1+p)\,\Theta+2p\,\qp(d')\pmod{p^2}
\]
(Theorem~\ref{thm:N}). The Fermat quotient $\qp(d')$ of the denominator is not determined by the image of $P$ in $E(\Qp)$.

All digits are described by one identity. Put $P''=pnP$, $F=(d'^{p-1}-1)/p$ and let $z$ be the formal logarithm. Then
\[
n^2h_p(P)=\frac1{1-p}\Bigl[p^{-2}\log(1+p^2\Theta)+2\log(1+pF)\Bigr]+p^{-2}\cL(z(P'')),
\]
where $\cL(z)=\log\bigl(x\,\sigma_p(t)^2\bigr)$ is an \emph{even} power series in $z$ whose coefficients are invariants of $E$: $\mu_2=-c$, $\mu_4=-b_4/12$, $\mu_6=-b_2b_4/360-b_6/30$, with $c=(b_2-E_2(E,\omega))/12$ (Theorem~\ref{thm:Nstar}, Proposition~\ref{prop:mu}). The congruences for the individual digits are truncations of this identity. In particular the normalisation of $\sigma_p$, through $E_2$, is visible only at the third digit, and the next local invariant $b_4$ only at the seventh. In rank one, $\min(v_p(\Reg_p),k)$ is a function of the layers $\gamma,\gamma_2,\dots,\gamma_k$ of the filtration and of $F\bmod p^{k-1}$, and of no smaller set of data (Theorem~\ref{thm:Gk}).

\subsection{A non-degeneracy}
Several of these statements need $v_p(\phi_p'(x(T)))=1$ at the points $T$ of order $p$. We prove that this holds for every $p\ge5$ and every ordinary curve over $\bar\F_p$ (Theorem~\ref{thm:ND}). The proof identifies a resultant of division polynomials modulo $p$ with a power of the discriminant times an explicit non-zero constant.

\subsection{An archimedean input}
A rational point in $E_k(\Qp)$ has naive height at least $2k\log p$. Hence, if some non-zero point of $\Lam\cap E_1(\Qp)$ has canonical height below $4\log p+\beta(E)$, where $\beta(E)$ is a lower bound for the local heights away from $p$, then $(\chi,\gamma)$ is surjective and, in rank two, $\Reg_p$ is a unit (Proposition~\ref{prop:canonical}). In our computations on curves of rank two this condition accounts, in expectation, for the whole observed frequency of non-unit regulators (Section~\ref{sec:arch}).

\subsection*{Conventions}
Results are proved on paper unless stated otherwise. The arithmetic steps are formalised in Lean~4 (Appendix~\ref{app:lean}); statements marked \emph{[computation]} were checked with PARI/GP on curves from Cremona's tables and are not proved. This paper and \cite{P1} are a revision of a draft that the author posted under the pen name Kiichi \cite{Draft}; Appendix~\ref{app:labels} lists the corresponding labels.

\section{Setting}\label{sec:setting}

We use the notation of \cite{P1}. For $P\in E(\Q)\smallsetminus\{O\}$ write $P=(\alpha/d^2,\beta/d^3)$ in lowest terms; $t=-x/y$; $E_k(\Qp)=\{v_p(t)\ge k\}$; $\log=\log_p$ is Iwasawa's logarithm. The height is $h_p(P)=2\log(\sigma_p(P)/d(P))$ for $P$ reducing to $O$ at $p$ and to non-singular points at the bad primes, extended quadratically; $\sigma_p\in t\Zp[[t]]$ is the Mazur--Tate sigma function \cite{MT91}, normalised by
\[
x(t)+c=-\frac{d}{\omega}\Bigl(\frac1\sigma\frac{d\sigma}{\omega}\Bigr),\qquad c=\frac{b_2-E_2(E,\omega)}{12},
\]
with $E_2(E,\omega)$ the value of the $p$-adic weight-two Eisenstein series. Division polynomials satisfy $x(mQ)=\phi_m/\psi_m^2$, $y(mQ)=\omega_m/\psi_m^3$; modulo $p$, $\phi_p\equiv f(x^p)$ and $\psi_p\equiv g(x^p)$.

For $u\in\Zp^\times$, $\qp(u)=(u^{p-1}-1)/p\bmod p$, and if $\qp(u)=0$ we put
\[
Q_2(u)=\frac{u^{p-1}-1}{p^2}\in\Zp,\qquad q^{(2)}(u)=Q_2(u)\bmod p .
\]

Throughout, unless stated otherwise:
\begin{itemize}
\item[(H)] $p\ge5$ is of good ordinary reduction, $\#\Et(\Fp)=p$, $p\nmid n$, ($\iota$) holds, and $E(\Qp)[p]\ne0$.
\end{itemize}
We write $D=E(\Qp)[p]$, fix $T_0\in D$ whose reduction $G$ generates $\Et(\Fp)$, and let $\chi$ be the discrete logarithm of the reduction with respect to $G$. For a class $\bar x\in\Et(\Fp)\smallsetminus\{O\}$ let $T(\bar x)\in D$ be the unique point above it.

\subsection{Results from \cite{P1}}\label{sec:recall}
We use: (a) for $P$ reducing to $O$ and satisfying the condition at bad primes, $h_p(P)\in2p\Zp$; (b) the decomposition $h_p(P)=2\log(\sigma_p(t)/t)+2\log(\alpha/\beta)$; (c) for $P'$ with reduction of order $p$, $v_p(t(pP'))\ge2$ in the split case; (d) $\lambda(\bar x)=\frac12\qp(\phi_p(x))=0$ for all lifts $x$ of all $\bar x\ne O$ in the split case, so $\phi_p(x)^{p-1}\in1+p^2\Zp$; (e) $\Mb$ vanishes on $\ker\chi\times\ker\chi$, and $h_p(e_1)\equiv2\lambda^{(2)}/n^2\pmod p$ for $\chi(e_1)\ne0$, where $\lambda^{(2)}=\frac12q^{(2)}(\phi_p(x(ne_1)))$; (f) for $p\nmid m$ and $Q\in E(\Zp)$ with non-zero reduction, $q^{(2)}(\phi_p(x(mQ)))=m^2q^{(2)}(\phi_p(x(Q)))$, proved as \cite[Prop.~6.1]{P1} with $\delta\in p^4\Zp$. We also use the formal-part bound $v_p(\log(\sigma_p(t)/t))\ge\min(e+v_p(a_1),2e)$, $e=v_p(t)$.

\section{The height matrix modulo \texorpdfstring{$p$}{p}}\label{sec:K}

\subsection{Two homomorphisms}
By (c) of \S\ref{sec:recall}, $D\cap E_1(\Qp)=0$, and $D\to\Et(\Fp)$ is an isomorphism. Hence $E(\Qp)=D\oplus E_1(\Qp)$. For $Q\in E(\Zp)$ write $Q=\chi(Q)T_0+R$ with $R\in E_1(\Qp)$ and put
\[
\gamma(Q)=t(R)/p\bmod p .
\]
Then $\chi,\gamma$ are homomorphisms, $\gamma(T_0)=0$, and $\ker\chi\cap\ker\gamma=E_2(\Qp)$. More generally, write $t(R)/p=\gamma+p\gamma_2+p^2\gamma_3+\cdots$; we call $\gamma_k$ the $k$-th layer.

\begin{lemma}\label{lem:firstdigit}
Let $P\in E(\Zp)$ have non-zero reduction and $R\in E_1(\Qp)$ with $t=t(R)$. Then $x(P+R)\in\Zp[[t]]$ and
\[
x(P+R)=x(P)+c_1(P)\,t+O(t^2),\qquad c_1(P)=2y(P)+a_1x(P)+a_3\in\Zp^\times .
\]
\end{lemma}

\begin{proof}
Translation by $P$ is an automorphism of the smooth model over $\Zp$ taking $O$ to a point with non-zero reduction, so $x\circ\tau_P\in\Zp[[t]]$. The invariant differential $\omega=dx/(2y+a_1x+a_3)=(1+a_1t+\cdots)dt$ is translation-invariant, so the coefficient of $t$ is $c_1(P)$. It is a unit because $\#\Et(\Fp)=p$ is odd, so $\Et(\Fp)$ has no point of order two.
\end{proof}

\subsection{The second Fermat quotient under a change of lift}
Put $G_\phi=\phi_p'/p$ and $H_\phi=G_\phi'$; both are polynomials with coefficients in $\Z[a_1,\dots,a_6]$. For a lift $x_0$ of $\bar x\ne O$ put $\Theta(x_0)=Q_2(\phi_p(x_0))$; by (d) this is defined.

\begin{proposition}\label{prop:lift}
For $k\in\Zp$, with $u=\phi_p(x_0)$,
\[
\Theta(x_0+pk)\equiv\Theta(x_0)+(p-1)\,k\,\frac{G_\phi(x_0)}u-p\,k^2\frac{H_\phi(x_0)}{2u}\pmod{p^2}.
\]
In particular $\lambda^{(2)}(x_0+pk)=\lambda^{(2)}(x_0)+k\,s(\bar x)$ with $s(\bar x)=-G_\phi(\bar x)/(2\phi_p(\bar x))\in\Fp$.
\end{proposition}

\begin{proof}
Taylor expansion gives $\phi_p(x_0+pk)\equiv u+p^2kG_\phi(x_0)+p^3k^2H_\phi(x_0)/2\pmod{p^4}$. Call the correction $\varepsilon$; then $p^2\mid\varepsilon$ and $(u+\varepsilon)^{p-1}\equiv u^{p-1}(1+(p-1)\varepsilon/u)\pmod{p^4}$. Divide by $p^2$.
\end{proof}

By Theorem~\ref{thm:ND} below, $s(\bar x)\ne0$ for every $\bar x$; so $\lambda^{(2)}$, unlike $\lambda$, depends on the lift.

\subsection{Two local identities and the matrix}
For $\bar x\ne O$ put $\lambda_T(\bar x)=\lambda^{(2)}(x(T(\bar x)))$, and let $B_0=s(G)\,c_1(T_0)$ and $\lambda_{T,1}=\lambda_T(G)$.

\begin{proposition}\label{prop:local}
For every $\bar x\in\Et(\Fp)\smallsetminus\{O\}$, with $m=\chi(\bar x)$,
\[
s(\bar x)\,c_1(T(\bar x))=B_0\,m,\qquad \lambda_T(\bar x)=\lambda_{T,1}\,m^2 .
\]
\end{proposition}

\begin{proof}
Let $L_2(Q)=\lambda^{(2)}(x(Q))$ for $Q\in E(\Zp)$ with non-zero reduction. By (f), $L_2(mQ)=m^2L_2(Q)$. By Lemma~\ref{lem:firstdigit}, $x(Q)\equiv x(T)+p\,c_1(T)\gamma(Q)\pmod{p^2}$ with $T=\chi(Q)T_0$, and by Proposition~\ref{prop:lift},
\[
L_2(Q)=\lambda_T(\bar x)+s(\bar x)\,c_1(T)\,\gamma(Q).
\]
On the coset $\chi=1$, $\gamma$ takes all values in $\Fp$. Replacing $Q$ by $mQ$ multiplies $\chi$ and $\gamma$ by $m$; comparing constant and linear terms in $\gamma(Q)$ of $L_2(mQ)=m^2L_2(Q)$ gives both identities.
\end{proof}

\begin{theorem}\label{thm:K}
Under \textup{(H)}, for all $P,R\in\Lam$,
\[
h_p(P)\equiv2\lambda_{T,1}\chi(P)^2+2B_0\,\chi(P)\gamma(P)\pmod p,
\]
\[
\Mb(P,R)=2\lambda_{T,1}\chi(P)\chi(R)+B_0\bigl(\chi(P)\gamma(R)+\gamma(P)\chi(R)\bigr).
\]
The constants depend only on $E/\Qp$ and $T_0$; the Tamagawa product $n$ does not appear.
\end{theorem}

\begin{proof}
For $\chi(P)\ne0$, (e) of \S\ref{sec:recall} gives $h_p(P)\equiv2L_2(nP)/n^2$. By the proof of Proposition~\ref{prop:local} and Proposition~\ref{prop:local} itself,
\[
L_2(nP)=\lambda_{T,1}n^2\chi(P)^2+B_0\,n\chi(P)\,n\gamma(P).
\]
For $\chi(P)=0$ both sides vanish modulo $p$ by (a). Polarise.
\end{proof}

\begin{corollary}\label{cor:K}
\begin{enumerate}[label=\textup{(\roman*)}]
\item $\operatorname{rank}\Mb=2$ if and only if $(\chi,\gamma)\colon\Lam\to\Fp^2$ is surjective.
\item If $r=1$, then $v_p(\Reg_p)\ge1$ if and only if $\gamma(e_1)/\chi(e_1)=-\lambda_{T,1}/B_0$.
\item Suppose $(\chi,\gamma)$ is surjective, let $\Lam_1=\Lam\cap E_2(\Qp)$ \textup{(}of rank $r-2$\textup{)} and $\widehat C$ the Gram matrix of $p^{-1}h_p$ on $\Lam_1$. Then $v_p(\Reg_p)=r-2$ if and only if $\det\widehat C\not\equiv0\pmod p$.
\end{enumerate}
\end{corollary}

\begin{proof}
(i) and (ii) are linear algebra over $\Fp$. For (iii) choose a basis $e_1$ with $\chi(e_1)=1$, $e_2\in\ker\chi$ with $\gamma(e_2)\ne0$, and a basis of $\Lam_1$. By Theorem~\ref{thm:K}, $M$ has the shape $\left(\begin{smallmatrix}a&b&0\\ b&pC_{22}&p*\\ 0&p*&p\widehat C\end{smallmatrix}\right)$ with $b$ a unit, and the Schur complement of the leading $2\times2$ block is $p\widehat C$ modulo $p^2$.
\end{proof}

Thus the deficit $r-2$ comes from the two layers $\chi$ and $\gamma$, and what remains in rank $r\ge3$ is a non-degeneracy of $p^{-1}h_p$ one layer deeper. [Computation: the exact formula $v_p(\Reg_p)=(r-2)+v_p(\det\widehat C)$ fails in a few boxes where $\det\widehat C$ is not a unit.]

\section{The second digit}\label{sec:second}

\begin{theorem}\label{thm:N}
Under \textup{(H)}, let $P\in\Lam$ with $\chi(P)\ne0$, put $P'=nP$, and write $x(P')=a/d'^2$ in lowest terms. Then
\[
n^2h_p(P)\equiv(1+p)\,\Theta(x(P'))+2p\,\qp(d')\pmod{p^2}.
\]
\end{theorem}

\begin{proof}
Take a model with $p\mid a_1$; put $P''=pP'$, $e=v_p(t(P''))\ge2$, $\phi=\phi_p(x(P'))$, $\psi=\psi_p(x(P'))$, $\omega=\omega_p(P')$. Four terms are needed modulo $p^4$.

(1) \emph{Formal part.} $2\log(\sigma_p(t)/t)\equiv a_1t(P'')\pmod{p^4}$, since the next term has valuation at least $2e\ge4$.

(2) \emph{Weierstrass equation.} Multiplying it at $P''$ by $\psi^6$, all terms of $\omega^2-\phi^3$ other than $-a_1\phi\omega\psi$ have valuation at least $2e\ge4$, so $\omega^2\equiv\phi^3-a_1\phi\omega\psi\pmod{p^4}$. Writing $\phi^{p-1}=1+p^2\Theta$ and $\omega^{p-1}=1+p^2B$, this gives $Q_2(\phi/\omega)\equiv-\Theta/2-a_1\omega\psi/(2p^2\phi^2)\pmod{p^2}$.

(3) \emph{Denominator.} As in \cite{P1}, $\alpha/\beta=\pm(\phi/\omega)d'^{-p^2}$, and $-p^2\log d'\equiv p^3\qp(d')\pmod{p^4}$.

(4) \emph{Logarithm.} For $u$ with $\qp(u)=0$, $\log u\equiv-p^2(1+p)Q_2(u)\pmod{p^4}$.

Adding, $h_p(P'')\equiv a_1t(P'')+(1+p)(p^2\Theta+a_1\omega\psi/\phi^2)+2p^3\qp(d')\pmod{p^4}$. Since $t(P'')=-\phi\psi/\omega$,
\[
a_1t(P'')+a_1\omega\psi/\phi^2=a_1\psi(\omega^2-\phi^3)/(\phi^2\omega),
\]
which has valuation at least $2v_p(a_1)+2e\ge6$ by (2), and $p\,a_1\omega\psi/\phi^2$ has valuation at least $4$. Hence $h_p(P'')\equiv(1+p)p^2\Theta+2p^3\qp(d')\pmod{p^4}$, and $h_p(P'')=p^2n^2h_p(P)$.
\end{proof}

Modulo $p$ this is (e). The term $\qp(d')$ depends on the global denominator, and two points with the same image in $E(\Qp)/E_k(\Qp)$ for every $k$ can have different $\qp(d')$. Thus no homomorphism $\Lam\cap E_2\to E_2/E_3$ determines the second digit. The local term, on the other hand, is governed by one $p$-adic number:

\begin{proposition}\label{prop:badlift}
Fix $\bar x\ne O$ and $T=T(\bar x)$. The map $\delta\mapsto\Theta(x(T)+p\delta)$ is a bijective isometry of $\Zp$. In particular it has a unique zero $\delta_*$, and $v_p(\Theta(x(T)+p\delta))=v_p(\delta-\delta_*)$.
\end{proposition}

\begin{proof}
By Proposition~\ref{prop:lift}, modulo $p$ this map is $\delta\mapsto\Theta(x(T))+2s\delta$ with $s\ne0$ (Theorem~\ref{thm:ND}), and its derivative is a unit everywhere. Hensel's lemma applied in each residue class gives the claim.
\end{proof}

\begin{corollary}\label{cor:Nprime}
Let $r=1$ and $v_p(\Reg_p)\ge1$, so that $p\mid\Theta(x(P'))$. Then $v_p(\Reg_p)\ge2$ if and only if $\Theta(x(P'))/p+2\qp(d')\equiv0\pmod p$.
\end{corollary}

\begin{remark}
Under $x\mapsto u^2x$ with $p\nmid u$, $\phi_p$ is multiplied by $u^{2p^2}$, so $\Theta\mapsto\Theta+2p\,\qp(u)\pmod{p^2}$, while $d'$ is unchanged if $\gcd(u,d')=1$. So Theorem~\ref{thm:N} refers to a model minimal at $p$; the term $\qp(d')$ is paired with the normalisation of the model.
\end{remark}

\section{A closed form for all digits}\label{sec:closed}

Keep the notation of Theorem~\ref{thm:N} and put
\[
\Theta=Q_2(\phi_p(x(P'))),\qquad F=\frac{d'^{p-1}-1}{p}\in\Z,\qquad z=\log_E(t(P'')),
\]
where $\log_E$ is the formal logarithm with $d\log_E=\omega$; so $v_p(z)=e$. Let $X=x(P'')$, $Y=y(P'')$, and
\[
R=\frac{a_2}X+\frac{a_4}{X^2}+\frac{a_6}{X^3}-\frac{a_1Y}{X^2}-\frac{a_3Y}{X^3},
\]
so that $\omega_p(P')^2=\phi_p(x(P'))^3(1+R)$.

\begin{theorem}\label{thm:Nstar}
Under the hypotheses of Theorem~\ref{thm:N},
\[
n^2h_p(P)=\frac1{1-p}\Bigl[p^{-2}\log(1+p^2\Theta)+2\log(1+pF)\Bigr]+p^{-2}\cL(z),
\]
where
\[
\cL(z)=\log\bigl(X\,\sigma_p(t(P''))^2\bigr)=2\log\frac{\sigma_p(t)}t-\log(1+R).
\]
\end{theorem}

\begin{proof}
From the proof of Theorem~\ref{thm:N}, $p^2n^2h_p(P)=h_p(P'')=2\log(\sigma_p(t)/t)+2\log(\phi/\omega)-2p^2\log d'$. Since $\omega^2=\phi^3(1+R)$, $2\log(\phi/\omega)=-\log\phi-\log(1+R)$. For a unit $u$, $\log u=(p-1)^{-1}\log(u^{p-1})$, so $-\log\phi-2p^2\log d'=-(p-1)^{-1}[\log(1+p^2\Theta)+2p^2\log(1+pF)]$. This gives the formula with the second expression for $\cL$. For the first, $t^2=X^2/Y^2=\psi^2/(\phi(1+R))=1/(X(1+R))$, so $X\sigma_p(t)^2=(\sigma_p(t)/t)^2/(1+R)$.
\end{proof}

\begin{proposition}\label{prop:mu}
As a power series in $z$, $\cL(z)=\sum_{m\ge1}\mu_{2m}z^{2m}$ has no odd terms, and
\[
\mu_2=-c,\qquad \mu_4=-\frac{b_4}{12},\qquad \mu_6=-\frac{b_2b_4}{360}-\frac{b_6}{30}.
\]
Moreover $\mu_{2m}\in\Z[1/(2m+2)!][b_2,b_4,b_6,E_2]$, and $E_2$ occurs only in $\mu_2$.
\end{proposition}

\begin{proof}
The defining equation of $\sigma_p$ is $(d/dz)^2\log\sigma_p=-(x+c)$, and \[x+b_2/12=\wp(z)=z^{-2}+\sum_{k\ge1}\pi_kz^{2k}\] with $\pi_1=c_4/240$, $\pi_2=c_6/6048$. Put $a=-b_2/12$. Then $(d/dz)^2\log(\sigma_p/z)=-(a+c)-\sum_k\pi_kz^{2k}$ is even, and integrating twice,
\[
2\log\frac{\sigma_p(z)}z=-(a+c)z^2-2\sum_{k\ge1}\frac{\pi_kz^{2k+2}}{(2k+1)(2k+2)} .
\]
Also $x(z)z^2=1+az^2+\pi_1z^4+\pi_2z^6+\cdots$ is even, so $\cL=\log(xz^2)+2\log(\sigma_p/z)$ is even. Comparing coefficients: at $z^2$, $a-(a+c)=-c$. At $z^4$, $\pi_1-a^2/2-\pi_1/6=(c_4-b_2^2)/288=-b_4/12$, using $c_4=b_2^2-24b_4$; here $E_2$ cancels. At $z^6$, $\pi_2-a\pi_1+a^3/3-\pi_2/15$; substituting $c_4$ and $c_6=-b_2^3+36b_2b_4-216b_6$, the terms in $b_2^3$ cancel and $\mu_6=-b_2b_4/360-b_6/30$. The denominators come from $\pi_k$, which has denominator dividing $(2k+2)!$, and from the factor $(2k+1)(2k+2)$.
\end{proof}

The denominators are not bounded independently of $m$: $\mu_6$ has a $5$, $\mu_8$ a $7$, and $\mu_{14}$ an $11$ and a $13$.

\begin{proposition}\label{prop:digits}
In Theorem~\ref{thm:Nstar}, the $j$-th term of $p^{-2}\log(1+p^2\Theta)$ has valuation at least $2j-2-v_p(j)$, the $j$-th term of $2\log(1+pF)$ has valuation at least $j-v_p(j)$, with equality in both when $\Theta$ and $F$ are units, and $p^{-2}\mu_{2m}z^{2m}$ has valuation $2me-2+v_p(\mu_{2m})$. A term of valuation $v$ first affects $n^2h_p(P)$ modulo $p^{v+1}$. For $e=2$, $\mu_2$ appears at the third digit, $\mu_4$ at the seventh, and $\mu_6$ at the tenth if $p=5$ and the eleventh if $p\ge7$.
\end{proposition}

\begin{proof}
Term-by-term valuation of the logarithm series, using $v_p(z)=e$; the last sentence uses $v_5(\mu_6)=-1$ and $v_p(\mu_6)\ge0$ for $p\ge7$. The series converges because $v_p(\mu_{2m})\ge-(2m+2)/(p-1)$.
\end{proof}

\begin{lemma}\label{lem:coords}
$z^2-1/x(z)=-(b_2/12)z^4+O(z^6)$ with only even powers. Hence $p^{-2}(z(P'')^2-1/x(P''))$ has valuation at least $4e-2\ge6$, and $1/x(P'')$ may replace $z(P'')^2$ through the sixth digit but not at the seventh.
\end{lemma}

\begin{proof}
Invert $x(z)z^2=1+az^2+\cdots$ with $a=-b_2/12$.
\end{proof}

Truncating Theorem~\ref{thm:Nstar} gives explicit congruences. With $c$ as above:

\begin{corollary}\label{cor:N34}
Under the hypotheses of Theorem~\ref{thm:N},
\begin{align*}
n^2h_p(P)&\equiv(1+p+p^2)\Theta+2p(1+p)F-p^2F^2-\tfrac{p^2}2\Theta^2-\frac{c}{p^2x(P'')}\pmod{p^3},\\
n^2h_p(P)&\equiv(1+p+p^2+p^3)\Theta+2p(1+p+p^2)F-p^2(1+p)F^2+\tfrac23p^3F^3\\
&\qquad-\tfrac{p^2}2(1+p)\Theta^2-\frac{c}{p^2x(P'')}\pmod{p^4}.
\end{align*}
\end{corollary}

\begin{proof}
Expand the three series of Theorem~\ref{thm:Nstar} using $(1-p)^{-1}\equiv1+p+p^2+p^3\pmod{p^4}$ and Proposition~\ref{prop:digits}: only $\mu_2z^2$ contributes from $\cL$, and by Lemma~\ref{lem:coords} it may be written $-c/x(P'')$. Every term on the right is invariant under $(x,y)\mapsto(x,y+sx)$.
\end{proof}

\begin{remark}
[Computation.] The truncations at the fifth, sixth and seventh digits, written out from Proposition~\ref{prop:digits}, hold in all $648$ boxes tested (conductor at most $12000$, $p\in\{5,7\}$, split, $e\ge2$). At the seventh digit $b_4$ enters through $\mu_4z^4$, and $1/x(P'')$ no longer serves in place of $z^2$. The identity of Theorem~\ref{thm:Nstar}, truncated, was checked to the tenth digit; $\cL$ needs no odd powers at any depth reached.
\end{remark}

\section{The minimal data at each digit}\label{sec:minimal}

Let $T=\chi(P)T_0$ and $t_1=t(nP-T)\in p\Zp$, with $t_1/p=\gamma+p\gamma_2+p^2\gamma_3+\cdots$.

\begin{proposition}\label{prop:layer}
$P''=p(nP-T)$, so $z(P'')=p\log_E(t_1)$ and $e=1+v_p(t_1)$. Moreover $t(P'')\equiv p^2\gamma\pmod{p^3}$, so $e=2$ if and only if $\gamma\ne0$, and $c/(p^4x(P''))\equiv c\gamma^2\pmod p$.
\end{proposition}

\begin{proof}
$T$ and $nP$ have the same reduction and $pT=O$. The formal logarithm satisfies $\log_E([p]t)=p\log_E(t)$. For the second statement, ordinary reduction gives height one, so the coefficients of $[p](t)=pt+\sum_{i\ge2}c_it^i$ are divisible by $p$ for $2\le i<p$; with $t_1=p\gamma+\cdots$ this gives $[p](t_1)\equiv p^2\gamma\pmod{p^3}$. Finally $1/x=t^2(1+O(t))$.
\end{proof}

So $\cL(z(P''))$ is a function of $t_1$, that is of the layers, and of the invariants $\mu_{2m}$; it adds no global datum.

\begin{theorem}\label{thm:Gprime}
For $t_1\in p\Zp$ and $j\ge1$, with $\Theta$ viewed as a function of $t_1$ through $x(T\oplus t_1)$,
\[
v_p\bigl(\Theta(t_1+p^j)-\Theta(t_1)\bigr)=j-1 .
\]
\end{theorem}

\begin{proof}
By Lemma~\ref{lem:firstdigit}, the $x$-coordinate moves by exactly $p^j$. Since $\phi_p\equiv f(x^p)$, all divided derivatives $\phi_p^{(k)}/k!$ with $1\le k\le p-1$ are divisible by $p$; so $\phi_p$ moves by $pG_\phi(x_0)\Delta+p\Delta^2B+\Delta^pC$ with $v_p(\Delta)=j$, which has valuation exactly $j+1$ because $G_\phi(x(T))$ is a unit (Theorem~\ref{thm:ND}). Raising the unit $\phi_p(x_0)$ to the power $p-1$ preserves this valuation, and dividing by $p^2$ gives $j-1$.
\end{proof}

\begin{proposition}\label{prop:rates}
Let $t_1'=t_1+p^j$ with $j\ge1$ and $z'=z(t_1')$. Then $v_p(z'-z)=1+j$ and
\[
v_p\bigl(p^{-2}(\cL(z')-\cL(z))\bigr)\ge j+1 .
\]
Hence, with $F$ fixed, the two sides of Theorem~\ref{thm:Nstar} at $t_1$ and $t_1'$ differ at valuation exactly $j-1$.
\end{proposition}

\begin{proof}
The Taylor expansion of $\log_E$ at $t_1$ has $k$-th term of valuation at least $kj-v_p(k)$, and the first term has valuation exactly $j$; multiply by $p$. For $m\ge1$, $z'^{2m}-z^{2m}$ has valuation at least $(1+j)+(2m-1)\min(e,1+j)\ge(1+j)+2(2m-1)$. With $v_p(\mu_{2m})\ge-(m+1)/2$, all terms of $p^{-2}\Delta\cL$ have valuation at least $j+1$. The $\Theta$-side moves at valuation $j-1$ by Theorem~\ref{thm:Gprime}, and the factor $(1-p)^{-1}$ and the bracket are units.
\end{proof}

\begin{theorem}\label{thm:Gk}
Assume \textup{(H)} and $r=1$. Fix $E$, $p$, $T_0$ and the local invariants. Then $\min(v_p(\Reg_p),k)$ is a function of
\[
(\chi,\gamma,\gamma_2,\dots,\gamma_k)\quad\text{and}\quad F\bmod p^{k-1},
\]
and of no smaller set of data: changing $\gamma_k$, or the $(k-1)$-st digit of $F$, alone can change $\min(v_p(\Reg_p),k)$.
\end{theorem}

\begin{proof}
\emph{Sufficiency.} By Theorem~\ref{thm:Nstar}, $v_p(\Reg_p)\ge k$ is a condition on $\Theta\bmod p^k$, on $F\bmod p^{k-1}$ and on $p^{-2}\cL$ modulo $p^k$. By Theorem~\ref{thm:Gprime}, $\Theta\bmod p^k$ depends on $t_1\bmod p^{k+1}$, and by Proposition~\ref{prop:layer} so does $p^{-2}\cL$ modulo $p^k$.

\emph{Necessity.} Changing $\gamma_k$ alone changes the right-hand side at valuation exactly $k-1$ by Proposition~\ref{prop:rates}, so at most one of the two values is divisible by $p^k$. Changing $F$ by $p^{k-2}$ changes $2\log(1+pF)$ at valuation exactly $k-1$, and nothing else moves.
\end{proof}

For $k=2,3,4$ the criteria are explicit. With $\Theta=p\Theta_1$, $S=(\Theta_1+2F)/p$ and $W=c/(p^4x(P''))\equiv c\gamma^2$:

\begin{corollary}\label{cor:criteria}
Let $r=1$.
\begin{enumerate}[label=\textup{(\roman*)}]
\item If $v_p(\Reg_p)\ge2$, then $v_p(\Reg_p)\ge3$ if and only if $S\equiv F^2+W\pmod p$.
\item If $v_p(\Reg_p)\ge3$ and $S-F^2-W=pU$, then $v_p(\Reg_p)\ge4$ if and only if $6U+6(S-F^2)+4F^3\equiv0\pmod p$.
\end{enumerate}
\end{corollary}

\begin{proof}
Substitute into Corollary~\ref{cor:N34}.
\end{proof}

\begin{remark}
On the rational points of a single curve the two kinds of data are related. Along $P+p^a\Lam$ the digit of $F$ at $p^{a-1}$ is determined by $\chi$, $n^2h_p(P)\bmod p$ and the homomorphism $Z\mapsto t(pnZ)/p^2\bmod p$, and $F\bmod p^k$ is determined by $n^2h_p(P)\bmod p^{k+1}$ and $t_1\bmod p^{k+2}$, by inverting Theorem~\ref{thm:Nstar}. Nevertheless both kinds of necessity in Theorem~\ref{thm:Gk} are realised by rational points: [computation] on the rank-two curve $4407a1$ at $p=5$, read point by point.
\end{remark}

\section{The non-degeneracy}\label{sec:ND}

We prove that $s(\bar x)\ne0$, equivalently $v_p(\phi_p'(x(T)))=1$, for all $p\ge5$. Work over $\bar\F_p$ with $\Et\colon y^2=u(x)=x^3+a_4x+a_6$ and $\Delta_0=4a_4^3+27a_6^2$, weights $2,4,6$ for $x,a_4,a_6$. Write $\psi_p\equiv\Psi(x^p)$ with $\deg\Psi=n_p=(p-1)/2$ and leading coefficient the Hasse invariant $A$; $h=\Psi/A$; $\phi_p\equiv\Phi(x^p)$; $w^{(p)}(X)=X^3+a_4^pX+a_6^p$. Then $h$ is the kernel polynomial of the Verschiebung $V$ and $\Phi/(A^2h^2)$ is its $x$-coordinate. Let $K=\psi_p'/p\bmod p$ and $S_1=(\phi_p-4uK^2)/\psi_p$; then the non-degeneracy at $\bar x$ is equivalent to $S_1(\bar x)\ne0$. Let
\[
\mathcal D=-6a_6\,\partial_{a_4}+\tfrac43a_4^2\,\partial_{a_6},
\]
which raises weight by $2$, kills $\Delta_0$, and acts on $q$-expansions of forms of weight divisible by $p-1$ as $-q\,d/dq$.

\begin{lemma}\label{lem:isog}
Let $\lambda\colon E'\to E$ be separable of odd degree with kernel polynomial $h$, $x\circ\lambda=P/h^2$, $y\circ\lambda=y'\Omega/h^3$, $\lambda^*\omega_E=c\,\omega_{E'}$. At every simple root $\xi$ of $h$ with $P(\xi)\ne0$,
\[
c^2P(\xi)=4u_{E'}(\xi)h'(\xi)^2,\qquad c^2P'(\xi)=2\bigl[u_{E'}(h')^2\bigr]'(\xi).
\]
\end{lemma}

\begin{proof}
Clearing denominators in $\omega_E=dx/2y$ gives $P'h-2Ph'=c\Omega$ and $u_{E'}\Omega^2=P^3+a_4Ph^4+a_6h^6$. Evaluate at $\xi$, and differentiate once.
\end{proof}

\begin{proposition}\label{prop:psi}
$A\cdot(\psi_p'/p)=u^{(p-1)/2}\,(d\Psi/dX)(x^p)$ in $\Fp[a_4,a_6][x]$.
\end{proposition}

\begin{proof}
It suffices to check this at ordinary points. Lift $\Et$ to $E$ over a complete discrete valuation ring with $v(p)=1$. Let $C=\widehat E[p]$ be the canonical subgroup. Its non-zero points have $v(x)=-2/(p-1)$, so $\psi_p=\lambda MN$ with $M=p\cdot(\text{kernel polynomial of }C)$, $M\equiv A_0$ a constant modulo $p$, and $N$ the factor for the other points. Write $\pi\colon E\to E/C$ and $\hat\pi$ with $\hat\pi\circ\pi=[p]$, $x\circ\pi=P_\pi/M^2$, $y\circ\pi=y\,\Omega_\pi/M^3$ and $c_\pi=p\varpi$; then $x\circ\pi\equiv x^p$ and $\Omega_\pi/M^3\equiv u^{n_p}$ modulo $p$. Up to a unit, $N=M^{2n_p}h^*(x\circ\pi)$ with $h^*\equiv h$. Differentiating and using $2n_p+1=p$,
\[
K=\lambda M^{2n_p}M'h^*(x\circ\pi)+\lambda\varpi M^{2n_p-2}(h^*)'(x\circ\pi)\,\Omega_\pi,
\]
Modulo $p$ the first term vanishes since $M'\equiv0$, and the second is $\kappa\,u^{n_p}h'(x^p)$ with $\kappa\in\F^\times$; comparing leading coefficients gives $\kappa=1$.
\end{proof}

\begin{proposition}\label{prop:S1}
$\phi_p\equiv4u(\psi_p'/p)^2+\mathcal D(\psi_p^2)\pmod p$; equivalently $S_1=2\mathcal D\psi_p$.
\end{proposition}

\begin{proof}
By Proposition~\ref{prop:psi} the claim is $\Phi-4w^{(p)}(h')^2=2A(\mathcal DA)h^2+2A^2h\,\mathcal Dh$. By Lemma~\ref{lem:isog} for $V$ the left side vanishes on the roots of $h$ and has degree at most $p-1$, so the identity amounts to
\[
\textup{(R)}\ \bigl[w^{(p)}(h')^2\bigr]'+A^2h'\,\mathcal Dh\equiv0\pmod h,\qquad \textup{(L)}\ [X^{p-1}]\bigl(\Phi-4w^{(p)}(h')^2\bigr)=2A\,\mathcal DA .
\]
Both sides are weight-homogeneous, and $q$-expansion is injective on each weight \cite{Ser73,SwD73}, so it suffices to check them on the Tate curve, where $A\equiv1$ and $V$ is the isogeny $E_{q^p}\to E_q$. With roots $\xi_i=X_Q(q^i)$, $\eta_i=Y_Q(q^i)$, $Q=q^p$: (a) $q\,d\xi_i/dq=2i\eta_i$, since $p\,Q\,\partial_Q=0$; (b) $H(W)=h(X_Q(W))=C'W^{n_p}\theta_q(W)\theta_Q(W)^{-p}$ by the product formula for $x_Q(W_1)-x_Q(W_2)$; (c) the logarithmic derivative of $H$ at $W=q^i$ has constant term $n_p+1-i$, which gives $\eta_i\,D^2H/DH=-2i\eta_i=\mathcal D\xi_i$, that is (R). For (L), the constant terms compare as $\phi_1+6\sigma$ with $\phi_1=-2\sigma$ and $\sigma=\sum\xi_i=-E_2/24$, while $2A\,\mathcal DA\equiv-E_2/6$.
\end{proof}

\begin{theorem}\label{thm:ND}
For every prime $p\ge5$ there is $c_S\in\Fp^\times$ with
\[
\mathrm{Res}_x(S_1,\psi_p\bmod p)=c_S\,\Delta_0^{p^2(p^2-1)/24}\quad\text{in }\Fp[a_4,a_6].
\]
Hence $v_p(\phi_p'(x(T)))=1$ for every ordinary curve and every $T$ of order $p$, and $s(\bar x)\ne0$.
\end{theorem}

\begin{proof}
By Proposition~\ref{prop:S1}, \[\mathrm{Res}_x(S_1,\psi_p)=\bigl(2^{n_p}\mathrm{Res}_X(\mathcal D\Psi,\Psi)\bigr)^p,\] and $\mathcal D\Psi(\xi_i)=-A(\mathcal D\xi_i)\prod_{j\ne i}(\xi_i-\xi_j)$. On the Tate curve, $\mathrm{ord}_q(\mathcal D\xi_i)=i$ and $\mathrm{ord}_q(\xi_i-\xi_j)=\min(i,j)$, so the resultant has $q$-order $M=p(p^2-1)/24$, which is one twelfth of its weight. As $\Delta_0$ is the only weight-homogeneous irreducible factor vanishing at the cusp, the resultant is $\gamma_p\Delta_0^M$ with $\gamma_p\ne0$; the leading coefficients give $\gamma_p=(-1)^{n_p(n_p+1)/2}\,n_p!\,(-16)^M$.
\end{proof}

[Computation: before this proof, the non-degeneracy had been verified for all curves over $\Fp$ with $p\le41$.] With $m_K=(p^2-1)(p^2-3)/24$ one also finds $\mathrm{Res}_x(K,\psi_p\bmod p)=\varepsilon_p2^{m_K}\Delta_0^{m_K}$ with $\varepsilon_p=-1$ exactly when $p\equiv3\pmod8$.

\section{An archimedean threshold}\label{sec:arch}

For $Q\in E(\Q)\smallsetminus\{O\}$ let $h(x(Q))=\log\max(|\alpha(Q)|,d(Q)^2)$.

\begin{lemma}\label{lem:depthheight}
If $p$ is a prime of good reduction and $Q\in E_k(\Qp)$, $k\ge1$, then $p^k\mid d(Q)$, so $h(x(Q))\ge2k\log p$.
\end{lemma}

\begin{proof}
$Q\in E_k(\Qp)$ means $v_p(x(Q))\le-2k$.
\end{proof}

\begin{proposition}\label{prop:naive}
Assume \textup{(H)} and that $E(\Q)$ is torsion-free, and put \[h_{\min}=\min\{h(x(Q)):0\ne Q\in\Lam\cap E_1(\Qp)\}.\] If $h_{\min}<4\log p$, then $(\chi,\gamma)$ is surjective; in rank two $\Reg_p$ is then a unit.
\end{proposition}

\begin{proof}
If $(\chi,\gamma)$ is not surjective, its image is the graph of $\gamma=c\chi$, so $\gamma$ vanishes on $\ker\chi=\Lam\cap E_1(\Qp)$, that is, $\Lam\cap E_1(\Qp)\subseteq E_2(\Qp)$. Lemma~\ref{lem:depthheight} with $k=2$ gives $h_{\min}\ge4\log p$. Then apply Corollary~\ref{cor:K}(i) and \cite[Cor.~6.6]{P1}.
\end{proof}

Write the canonical height (normalised so that $\hat h\approx h(x)$) as $\hat h(Q)=2v_p(d(Q))\log p+2\log c'+\lambda'_\infty(Q)+\sum_{\ell}\mathrm{corr}_\ell(Q)\log\ell$, with $d(Q)=p^{v_p(d)}c'$ and $\mathrm{corr}_\ell\le0$ the component corrections \cite{SilH}. Let $m_\infty=\min_{E(\mathbb R)}\lambda'_\infty$, let $m_\ell$ be the minimum of $\mathrm{corr}_\ell$ over the components, and put $\beta(E)=m_\infty+\sum_\ell m_\ell\log\ell$.

\begin{proposition}\label{prop:canonical}
Under the hypotheses of Proposition~\ref{prop:naive}, if
\[
\hat h_{\min}=\min\{\hat h(Q):0\ne Q\in\Lam\cap E_1(\Qp)\}<4\log p+\beta(E),
\]
then $(\chi,\gamma)$ is surjective, and in rank two $\Reg_p$ is a unit.
\end{proposition}

\begin{proof}
If $(\chi,\gamma)$ is not surjective, every non-zero $Q\in\Lam\cap E_1(\Qp)$ has $v_p(d(Q))\ge2$, and the remaining terms are at least $\beta(E)$.
\end{proof}

\begin{lemma}\label{lem:imageE2}
Under \textup{(H)}, $E(\Qp)/E_2(\Qp)\cong(\Z/p)^2$, and $(\chi,\gamma)$ fails to be surjective if and only if the image of $\Lam$ in $E(\Qp)/E_2(\Qp)$ has order $p$. This depends only on the images of generators of any subgroup of $\Lam$ of index prime to $p$.
\end{lemma}

\begin{proof}
$E(\Qp)=D\oplus E_1(\Qp)$ and $E_1/E_2\cong\Z/p$. The image of $\Lam$ has order $p\cdot\#(\text{image of }\Lam\cap E_1)$, and the second factor is trivial exactly when $\Lam\cap E_1\subseteq E_2$.
\end{proof}

[Computation.] Over $3{,}056$ split boxes of rank two from Cremona's tables, Proposition~\ref{prop:naive} accounts for about $68\%$ of the observed deficit in the frequency of $\det(\chi,\gamma)\equiv0$ relative to $1/p$, and Proposition~\ref{prop:canonical} for $100.6\%$ in expectation; on a disjoint set of $5{,}572$ boxes (conductor between $200{,}000$ and $500{,}000$) the figures are $55.8\%$ and $99.7\%$. No box with $\det(\chi,\gamma)\equiv0$ lies below the canonical threshold, in these sets or in two further sets of $5{,}782$ and $7{,}861$ boxes outside Cremona's tables. Above the threshold the frequency is $1/p$ within the noise on the first two sets; on the last two the non-surjective boxes cluster just above the threshold, for a reason we do not know. The constant $m_\infty$ was computed numerically.

\section{Open questions}\label{sec:open}

\begin{enumerate}[label=(\arabic*)]
\item Is $\det\widehat C$ in Corollary~\ref{cor:K}(iii) governed by any local datum? It is a non-degeneracy of the same kind as Schneider's conjecture.
\item Find a sharp lower bound for $v_p(\mu_{2m})$, which would determine the digit of every local invariant.
\item Explain the frequency $1/p$ of $\det(\chi,\gamma)\equiv0$ above the canonical threshold, and the clustering just above it.
\item [Computation] For every anomalous curve over $\Fp$ with $p\le157$, the slope $\rho$ of \cite[App.~A]{P1} satisfies $(\rho/p)=(-3/p)$. Is there a closed form?
\end{enumerate}

\appendix

\section{Lean formalisation}\label{app:lean}

The Lean project for this paper has $25$ files, each importing only Mathlib (Lean~4 \texttt{v4.33.1}, Mathlib \texttt{0df444a3}), with no \texttt{sorry} or \texttt{native\_decide}; all \texttt{\#print axioms} lines are subsets of \{\texttt{propext}, \texttt{Classical.choice}, \texttt{Quot.sound}\}. Sources and axiom logs: \url{https://github.com/yoshikazuhosomura-sys/padic-height-anomalous}, doi:10.5281/zenodo.23856616. Lean checks the arithmetic steps; the $p$-adic sigma function, division polynomials, formal groups, the Tate curve and local heights enter as hypotheses. File $k$ is \texttt{PadicHeightValuation}$k$\texttt{.lean}.

\begin{center}\small
\begin{longtable}{p{0.28\linewidth}p{0.54\linewidth}p{0.08\linewidth}}
\toprule
Result & Lean declarations & File\\
\midrule
Prop.~\ref{prop:lift} & \lean{eval_shift_of_deriv_eq}, \lean{fermatQuot2_move}, \lean{lambda2_affine} & 11\\
Lemma~\ref{lem:firstdigit}, Thm.~\ref{thm:K}, Cor.~\ref{cor:K} & \lean{eval_shift_first_digit}, \lean{master}, \lean{master_vanish_on_ker_chi}, \lean{det_three_layer} & 12\\
Thm.~\ref{thm:N}, Prop.~\ref{prop:badlift} & \lean{a1_correction_cancels}, \lean{second_digit_criterion}, \lean{root_factor}, \lean{unique_bad_class} & 13\\
Cor.~\ref{cor:Nprime} & \lean{trichotomy}, \lean{sq_dvd_sub_crit} & 14\\
Cor.~\ref{cor:N34}, \ref{cor:criteria} & \lean{denominator_third_digit}, \lean{q2_quotient}, \lean{cube_dvd_iff_crit}, \lean{denominator_fourth_digit}, \lean{quartic_dvd_iff_crit}, \lean{quintichotomy} & 15, 16\\
Prop.~\ref{prop:layer} & \lean{p_mult_second_digit}, \lean{w_local}, \lean{layer_depth} & 16, 23\\
Thm.~\ref{thm:Gprime} & \lean{shift_depth_exact}, \lean{unit_pow_diff_exact}, \lean{theta_sensitivity_exact} & 17\\
Thm.~\ref{thm:ND} & \lean{resultant_weight}, \lean{nondegeneracy_summary}, \lean{nd_iff_S_ne_zero} & 18, 19\\
Prop.~\ref{prop:mu}, \ref{prop:digits} & \lean{mu2_eq}, \lean{mu4_eq}, \lean{mu6_eq}, \lean{mu6_no_b2_cube}, \lean{log_term_factor}, \lean{mu6_visible_at_ten} & 22, 23\\
Prop.~\ref{prop:rates}, Thm.~\ref{thm:Gk} & \lean{ledger_lags}, \lean{ledger_lags_any_m}, \lean{data_count}, \lean{layer_is_needed}, \lean{fermat_digits_needed} & 24, 25\\
Lemma~\ref{lem:depthheight}, Prop.~\ref{prop:naive} & \lean{height_ge_of_depth}, \lean{surjective_of_shallow_kernel_element} & 27\\
Prop.~\ref{prop:canonical} & \lean{surjective_of_canonical_lt_four_log_add_beta} & Can.\\
Remark after Cor.~\ref{cor:criteria} & \lean{translate_quadratic}, \lean{ledger_digit}, \lean{inverse_n3} & 29\\
\bottomrule
\end{longtable}
\end{center}

\section{Correspondence with the draft}\label{app:labels}

\begin{center}\small
\begin{tabular}{ll|ll}
\toprule
This paper & \cite{Draft} & This paper & \cite{Draft}\\
\midrule
Lemma~\ref{lem:firstdigit} & Lemma 7 & Prop.~\ref{prop:layer} & Props.~12, 19\\
Prop.~\ref{prop:lift} & Props.~9, R$''''$ & Thm.~\ref{thm:Gprime} & Thm.~G$'$\\
Prop.~\ref{prop:local} & Prop.~8 & Prop.~\ref{prop:rates} & Prop.~20\\
Thm.~\ref{thm:K}, Cor.~\ref{cor:K} & Thm.~K, Cor.~K$'$ & Thm.~\ref{thm:Gk} & Thms.~G$''$, G$^{(k)}$\\
Thm.~\ref{thm:N} & Thm.~N & Cor.~\ref{cor:criteria} & Cors.~N$''$, N$'''$\\
Prop.~\ref{prop:badlift} & Prop.~10 & Lemma~\ref{lem:isog} & Lemma 13\\
Cor.~\ref{cor:Nprime} & Cor.~N$'$ & Props.~\ref{prop:psi}, \ref{prop:S1} & Props.~14, 15\\
Thm.~\ref{thm:Nstar} & Thm.~N$^*$ & Thm.~\ref{thm:ND} & Thm.~ND, Prop.~16\\
Prop.~\ref{prop:mu}, \ref{prop:digits} & Props.~17, 18 & Lemma~\ref{lem:depthheight} & Lemma 23\\
Lemma~\ref{lem:coords} & Lemma 21 & Props.~\ref{prop:naive}, \ref{prop:canonical} & Props.~24, 24$^c$\\
Cor.~\ref{cor:N34} & Cors.~N$^{(3)}$, N$^{(4)}$ & Lemma~\ref{lem:imageE2} & Lemma 25\\
\bottomrule
\end{tabular}
\end{center}

\section*{Use of AI}
This research was carried out with the assistance of Claude Code (Anthropic). The constructions, proofs, computations and the Lean formalisation were developed with the AI system under the author's direction, and separate sessions of the system re-checked the Lean development and the computations. The author reviewed the arguments, checked the Lean development and its axiom logs, and takes full responsibility for the content.

\begin{thebibliography}{SwD73}
\bibitem[Draft]{Draft} Kiichi (Y.~Hosomura), \emph{Integrality of the $p$-adic height pairing at an anomalous prime is equivalent to a point of order $p$ in $E(\Qp)$}, draft v15, computo ergo sum, 2026.
\bibitem[MT91]{MT91} B.~Mazur, J.~Tate, \emph{The $p$-adic sigma function}, Duke Math. J. 62 (1991), 663--688.
\bibitem[P1]{P1} Y.~Hosomura, \emph{The $p$-adic height pairing at anomalous primes and points of order $p$ over $\Qp$}, 2026.
\bibitem[Ser73]{Ser73} J.-P. Serre, \emph{Formes modulaires et fonctions zêta $p$-adiques}, Lecture Notes in Math.~350, Springer, 1973, 191--268.
\bibitem[SilH]{SilH} J.~H. Silverman, \emph{Computing heights on elliptic curves}, Math. Comp. 51 (1988), 339--358.
\bibitem[SwD73]{SwD73} H.~P.~F. Swinnerton-Dyer, \emph{On $\ell$-adic representations and congruences for coefficients of modular forms}, Lecture Notes in Math.~350, Springer, 1973, 1--55.
\end{thebibliography}
\end{document}
